\section{Best Time to Buy and Sell Stock with Transaction Fee}
\label{sec:dp-stock-6}

\problemstatement{Given $\textit{price}$ and a fixed $\textit{fee}$,
find the maximum profit from unlimited transactions, where every
\textbf{completed} transaction (one buy plus one sell) costs
$\textit{fee}$.}

\begin{patternbox}
This closing problem of the chapter makes one final, equally small
modification to Section~\ref{sec:dp-stocks-intro}'s base template: the
sell clause now subtracts $\textit{fee}$ from the proceeds, charging the
cost of the completed round-trip at the moment it completes.
\end{patternbox}

\subsection*{Deriving the recurrence}

\[
  f(i, \textit{holding}) =
  \begin{cases}
    \max\big(-\textit{price}[i] + f(i{+}1, 1),\ f(i{+}1, 0)\big) &
    \text{not holding}, \\
    \max\big(\textit{price}[i] - \textit{fee} + f(i{+}1, 0),\
    f(i{+}1, 1)\big) & \text{holding.}
  \end{cases}
\]
Only the sell clause changed: $\textit{price}[i] \to \textit{price}[i]
- \textit{fee}$.

\subsection*{Approach 1: Recursion}

\begin{lstlisting}
#include <bits/stdc++.h>
using namespace std;

int stock6Recursive(int i, int holding, vector<int>& price, int fee) {
    int n = price.size();
    if (i == n) return 0;

    if (holding) {
        return max(price[i] - fee + stock6Recursive(i + 1, 0, price, fee),
                   stock6Recursive(i + 1, 1, price, fee));
    }
    return max(-price[i] + stock6Recursive(i + 1, 1, price, fee),
               stock6Recursive(i + 1, 0, price, fee));
}

int main() {
    vector<int> price = {1, 3, 2, 8, 4, 9};
    int fee = 2;
    cout << stock6Recursive(0, 0, price, fee) << endl; // 8
    return 0;
}
\end{lstlisting}

\begin{complexitybox}[Approach 1 Complexity]
\textbf{Time.} $O(2^n)$.
\textbf{Space.} $O(n)$ recursion stack.
\end{complexitybox}

\subsection*{Approach 2: Memoization}

\begin{lstlisting}
#include <bits/stdc++.h>
using namespace std;

int stock6Memo(int i, int holding, vector<int>& price, int fee,
               vector<vector<int>>& memo) {
    int n = price.size();
    if (i == n) return 0;
    if (memo[i][holding] != -1) return memo[i][holding];

    int result;
    if (holding) {
        result = max(price[i] - fee + stock6Memo(i + 1, 0, price, fee, memo),
                      stock6Memo(i + 1, 1, price, fee, memo));
    } else {
        result = max(-price[i] + stock6Memo(i + 1, 1, price, fee, memo),
                      stock6Memo(i + 1, 0, price, fee, memo));
    }
    return memo[i][holding] = result;
}

int main() {
    vector<int> price = {1, 3, 2, 8, 4, 9};
    int n = price.size(), fee = 2;
    vector<vector<int>> memo(n, vector<int>(2, -1));
    cout << stock6Memo(0, 0, price, fee, memo) << endl; // 8
    return 0;
}
\end{lstlisting}

\begin{complexitybox}[Approach 2 Complexity]
\textbf{Time.} $O(n)$.
\textbf{Space.} $O(n)$ memo $+$ $O(n)$ recursion stack.
\end{complexitybox}

\subsection*{Approach 3: Tabulation}

\begin{lstlisting}
#include <bits/stdc++.h>
using namespace std;

int stock6Tabulation(vector<int>& price, int fee) {
    int n = price.size();
    vector<vector<int>> dp(n + 1, vector<int>(2, 0));

    for (int i = n - 1; i >= 0; i--) {
        dp[i][1] = max(price[i] - fee + dp[i + 1][0], dp[i + 1][1]);
        dp[i][0] = max(-price[i] + dp[i + 1][1], dp[i + 1][0]);
    }
    return dp[0][0];
}

int main() {
    vector<int> price = {1, 3, 2, 8, 4, 9};
    cout << stock6Tabulation(price, 2) << endl; // 8
    return 0;
}
\end{lstlisting}

\subsection*{Dry run}

\dryrun{Take $\textit{price} = [1,3,2,8,4,9]$, $\textit{fee}=2$.}

The optimal strategy: buy at $1$ (day $0$), sell at $8$ (day $3$,
proceeds $8-2=6$ after fee, profit $6-1=5$), buy at $4$ (day $4$), sell
at $9$ (day $5$, proceeds $9-2=7$, profit $7-4=3$): total $5+3=8$.
Notice the fee makes the intermediate sell at price $3$ (day $1$)
unattractive -- selling and immediately rebuying would pay the fee
twice for little gain, so the table correctly holds through day $1$ and
$2$ to wait for the larger rise to $8$.

\begin{complexitybox}[Approach 3 Complexity]
\textbf{Time.} $O(n)$.
\textbf{Space.} $O(n)$ table, no recursion stack.
\end{complexitybox}

\subsection*{Approach 4: Space Optimization}

\begin{lstlisting}
#include <bits/stdc++.h>
using namespace std;

int stock6SpaceOptimized(vector<int>& price, int fee) {
    int n = price.size();
    int nextHold = 0, nextNotHold = 0;

    for (int i = n - 1; i >= 0; i--) {
        int curHold = max(price[i] - fee + nextNotHold, nextHold);
        int curNotHold = max(-price[i] + nextHold, nextNotHold);
        nextHold = curHold;
        nextNotHold = curNotHold;
    }
    return nextNotHold;
}

int main() {
    vector<int> price = {1, 3, 2, 8, 4, 9};
    cout << stock6SpaceOptimized(price, 2) << endl; // 8
    return 0;
}
\end{lstlisting}

\begin{complexitybox}[Approach 4 Complexity]
\textbf{Time.} $O(n)$.
\textbf{Space.} $O(1)$.
\end{complexitybox}

\begin{takeawaybox}
This chapter's six variants, viewed together, demonstrate exactly how
far a single well-chosen state (\textit{day}, \textit{holding}, plus an
optional transaction counter) can stretch: a transaction cap adds a
dimension; a cooldown changes a jump distance; a fee changes a single
arithmetic term. None of the six problems needed a genuinely new way of
thinking about the state -- each is the same holding-state machine,
read closely for which single detail the problem statement perturbed.
\end{takeawaybox}
